\section{BST Iterator}
\label{sec:bst-iterator}

\problemstatement{Design an iterator over a BST that supports
$\textit{next}()$ (returns the next smallest value in the tree, in
sorted order) and $\textit{hasNext}()$ (whether any values remain),
initialized with the tree's root. Across the iterator's whole lifetime,
$\textit{next}$ should average $O(1)$ time and the iterator should use
$O(h)$ space, not $O(n)$.}

\begin{patternbox}
\emph{``Iterator returning sorted values one at a time''} is exactly
Section~\ref{sec:bst-kth-smallest}'s iterative inorder traversal,
\textbf{paused and resumed} across separate method calls instead of run
to completion inside a single function. The keyword shift from
``compute the $k$-th value'' to ``design an iterator'' is a signal that
the same stack-based traversal state needs to persist as object state
between calls, rather than living entirely within one function's local
variables.
\end{patternbox}

\subsection*{Understanding the problem carefully}

If we simply ran a full inorder traversal up front and stored all $n$
values in a list, $\textit{next}()$ would be trivially $O(1)$, but
construction would cost $O(n)$ time \emph{and} space -- failing the
$O(h)$ space requirement. The fix is to keep \emph{only} the explicit
stack from Section~\ref{sec:bst-kth-smallest}'s iterative traversal alive
as the iterator's state, advancing it by exactly one step per
$\textit{next}()$ call instead of running it to completion all at once.

\begin{ideabox}[Persist the Traversal Stack as Iterator State]
On construction: push the entire leftward chain from the root onto the
stack (exactly as in Section~\ref{sec:bst-kth-smallest}'s setup).
$\textit{hasNext}()$: simply check whether the stack is non-empty.
$\textit{next}()$: pop the top of the stack (this is the next value in
sorted order); then push the leftward chain starting from the popped
node's right child (preparing the stack so that the \emph{following}
call to \textit{next} will correctly find the next value after this
one); return the popped value.
\end{ideabox}

\subsection*{Why pushing only the new leftward chain after each pop keeps amortized cost low}

Each node in the tree is pushed onto the stack exactly once during the
iterator's entire lifetime (either during construction, for nodes on the
initial leftward chain, or during some later $\textit{next}()$ call,
when it becomes the leftward-chain push triggered by its parent being
popped), and each node is popped exactly once. So the \emph{total} work
across all calls to $\textit{next}()$ over the iterator's lifetime is
$O(n)$ (each node contributes $O(1)$ amortized push/pop work) -- which,
spread over the $n$ calls needed to exhaust the tree, gives $O(1)$
\emph{average} time per call, even though any single call might cost up
to $O(h)$ if it happens to push a long leftward chain.

\subsection*{Algorithm}

\begin{enumerate}
  \item \textsc{Constructor}(\textit{root}): push the leftward chain
        from \textit{root} onto the stack.
  \item \textsc{HasNext}(): return whether the stack is non-empty.
  \item \textsc{Next}(): pop $\textit{node}$ from the stack; push the
        leftward chain starting from $\textit{node}.\textit{right}$;
        return $\textit{node}.\textit{val}$.
\end{enumerate}

\subsection*{Dry run}

\dryrun{Take the BST with root $7$, left child $3$, right child $15$
(left child $9$, right child $20$). Call \textit{next} three times.}

\begin{itemize}
  \item Construct: push leftward chain from $7$: push $7$, push $3$ (no
        left child). Stack (top to bottom) $=[3,7]$.
  \item \textit{next}(): pop $3$. Push leftward chain from $3$'s right
        (null) -- nothing to push. Stack $=[7]$. Return $3$.
  \item \textit{next}(): pop $7$. Push leftward chain from $7$'s right,
        $15$: push $15$, then push $15$'s left, $9$ (no left child).
        Stack $=[9,15]$. Return $7$.
  \item \textit{next}(): pop $9$. Push leftward chain from $9$'s right
        (null) -- nothing. Stack $=[15]$. Return $9$.
\end{itemize}

Output sequence so far: $3,7,9$ -- correctly sorted order.

\subsection*{C++ implementation}

\begin{lstlisting}
#include <bits/stdc++.h>
using namespace std;

struct TreeNode {
    int val;
    TreeNode* left;
    TreeNode* right;
    TreeNode(int v) : val(v), left(nullptr), right(nullptr) {}
};

class BSTIterator {
    stack<TreeNode*> st;

    void pushLeftChain(TreeNode* node) {
        while (node) {
            st.push(node);
            node = node->left;
        }
    }

public:
    BSTIterator(TreeNode* root) {
        pushLeftChain(root);
    }

    bool hasNext() {
        return !st.empty();
    }

    int next() {
        TreeNode* node = st.top();
        st.pop();
        pushLeftChain(node->right);
        return node->val;
    }
};

int main() {
    TreeNode* root = new TreeNode(7);
    root->left = new TreeNode(3);
    root->right = new TreeNode(15);
    root->right->left = new TreeNode(9);
    root->right->right = new TreeNode(20);

    BSTIterator it(root);
    while (it.hasNext()) cout << it.next() << " ";
    cout << endl;
    return 0;
}
\end{lstlisting}

\begin{complexitybox}
\textbf{Time Complexity.} $O(1)$ average per $\textit{next}()$ call
(amortized over the iterator's lifetime), $O(h)$ worst case for any
single call; $\textit{hasNext}()$ is always $O(1)$.
\textbf{Space Complexity.} $O(h)$ for the stack, never $O(n)$.
\end{complexitybox}

\begin{takeawaybox}
Converting a one-shot traversal into a resumable iterator is purely a
matter of moving the traversal's local state (here, the explicit stack)
into object-level state that persists between calls. The traversal logic
itself -- push the leftward chain, pop, push the popped node's right
subtree's leftward chain -- is completely unchanged from
Section~\ref{sec:bst-kth-smallest}.
\end{takeawaybox}
